% Part: normal-modal-logic
% Chapter: filtrations
% Section: preliminaries

\documentclass[../../../include/open-logic-section]{subfiles}

\begin{document}

\olfileid[hi]{nml}{fil}{pre}

\olsection{प्रारंभिक बातें}

निस्यंदन हमें यह दिखाकर अपने मोडल तर्कशास्त्रीय तंत्रों की निर्णेयता
स्थापित करने देते हैं कि उनमें \emph{परिमित प्रतिरूप गुण} है, अर्थात् किसी
प्रतिरूप में सत्य (असत्य) हर !!{formula} किसी \emph{परिमित} प्रतिरूप में भी
सत्य (असत्य) है। निस्यंदन, !!{formula}ों के ऐसे समुच्चयों के सापेक्ष
परिभाषित होते हैं जो उपसूत्रों के अधीन संवृत हैं।

\begin{defn}\ollabel{defn:modallyclosed}
  !!{formula}ों का कोई समुच्चय $\Gamma$ \emph{उपसूत्रों के अधीन संवृत} है
  यदि उसमें $\Gamma$ के प्रत्येक !!a{formula} का प्रत्येक उपसूत्र हो। इसके
  अतिरिक्त, $\Gamma$ \emph{मोडलतः संवृत} है यदि वह उपसूत्रों के अधीन
  संवृत हो और $!A \in \Gamma$ से $\Box!A, \Diamond!A \in \Gamma$ भी
  निकलता हो।
\end{defn}

उदाहरणार्थ, कोई !!a{formula}~$!A$ दिए होने पर उसके सभी उप-!!{formula}ों का
समुच्चय उप-!!{formula}ों के अधीन संवृत है। जब हम $!A$ के उप-!!{formula}ों के
समुच्चय के माध्यम से किसी प्रतिरूप का निस्यंदन परिभाषित करते हैं, तो उसमें
हमारा अपेक्षित गुण होगा: वह $!A$ को तभी सत्य (असत्य) बनाता है जब मूल
प्रतिरूप ऐसा करता है।

$\Gamma$ के माध्यम से $\mModel{M}$ के निस्यंदन के जगतों का समुच्चय
निम्न तुल्यता संबंध के सभी तुल्यता-वर्गों के समुच्चय के रूप में परिभाषित
होता है।

\begin{defn}
मान लें $\mModel{M} =\tuple{W, R, V}$ और $\Gamma$ उप-!!{formula}ों के अधीन
संवृत है। $W$ पर संबंध $\equiv$ को उन किन्हीं दो जगतों के बीच सत्य
परिभाषित करें जो $\Gamma$ के उन्हीं !!{formula}ों को सत्य बनाते हैं,
अर्थात्:
\[
u \equiv v \quad \text{तभी जब } \quad 
\forall !A \in \Gamma : \mSat{M}{!A}[u] \Leftrightarrow \mSat{M}{!A}[v].
\]
किसी जगत~$w$ का तुल्यता-वर्ग~$[w]_\equiv$, या संक्षेप में $[w]$, उन सभी
जगतों का समुच्चय है जो $w$ के $\equiv$-तुल्य हैं:
\[
[w] = \Setabs{v}{v \equiv w}.
\]
\end{defn}

\begin{prop}
  $\mModel{M}$ और $\Gamma$ दिए होने पर, ऊपर परिभाषित $\equiv$ एक
  तुल्यता संबंध है, अर्थात् वह स्वपरावर्ती, सममित और संक्रमणशील है।
\end{prop}

\begin{proof}
  संबंध $\equiv$ स्वपरावर्ती है, क्योंकि $w$, $\Gamma$ के ठीक उन्हीं
  !!{formula}ों को सत्य बनाता है जिन्हें वह स्वयं सत्य बनाता है। वह सममित
  है, क्योंकि यदि $u$, $\Gamma$ के उन्हीं !!{formula}ों को सत्य बनाता है
  जिन्हें $v$, तो $v$ और $u$ के लिए भी यही सत्य है। वह संक्रमणशील भी है,
  क्योंकि यदि $u$, $\Gamma$ के उन्हीं !!{formula}ों को सत्य बनाता है जिन्हें
  $v$, और $v$ उन्हीं को जिन्हें $w$, तो $u$ भी उन्हीं !!{formula}ों को सत्य
  बनाता है जिन्हें $w$।
\end{proof}

हर तुल्यता संबंध की तरह संबंध $\equiv$, $W$ को \emph{विभाजनों} में
बाँटता है, अर्थात् $W$ के ऐसे उपसमुच्चयों में जो युग्मतः असंयुक्त हैं और
मिलकर पूरे $W$ को आच्छादित करते हैं। प्रत्येक $w \in W$ किसी एक विभाजन,
अर्थात् $[w]$, का !!a{element} है, क्योंकि $w \equiv w$। अतः विभाजन $[w]$
पूरे $W$ को आच्छादित करते हैं। वे युग्मतः असंयुक्त हैं, क्योंकि यदि $u
\in [w]$ और $u \in [v]$, तो $u \equiv w$ और $u \equiv v$; और सममिति
तथा संक्रमणशीलता से $w \equiv v$, अतः $[w] = [v]$।

\end{document}
